% Imported from algorithms/inference-optimization/README.md
\chapter{推理优化与瓶颈分析}\label{seedance-m10}
\moduleinfo{M10}{基础方法与机制解释}{\hyperref[src-seedance-1.0]{\texttt{seedance-1.0}}、\hyperref[src-seedance-1.5-pro]{\texttt{seedance-1.5-pro}}、\hyperref[src-rayflow]{\texttt{rayflow}}、\hyperref[src-adaspa]{\texttt{adaspa}}、\hyperref[src-seaweed-apt]{\texttt{seaweed-apt}}、\hyperref[src-longlive]{\texttt{longlive}}、\hyperref[src-wan21-code]{\texttt{wan21-code}}、\hyperref[src-diffsynth-code]{\texttt{diffsynth-code}}、\hyperref[src-longlive-code]{\texttt{longlive-code}}}

速度输出进入采样回路后，再看哪些计算重复、哪些近似可以减少调用。

采样器更快，用户是不是就等得更短？先把整次请求写出来。输入解码、提示改写、参考编码、去噪、精修、VAE 解码、音轨处理、上传和排队，都可能占时间。若真正慢的部分没有改，优化一个已经很快的环节，总耗时也不会明显变化。

\section{用端到端路径建立测量基线}\label{seedance-m10-ux7528ux7aefux5230ux7aefux8defux5f84ux5efaux7acbux6d4bux91cfux57faux7ebf}

因此先按一次请求的实际顺序把耗时拆开：

\begin{verbatim}
排队 → 输入处理 → 文本与参考编码 → 基础采样
     → 可选精修 → 视频/音频解码 → 编码混流 → 返回
\end{verbatim}

接着同时看墙钟与 GPU 内部时间。GPU 异步执行，计时边界不作同步，某段就可能显得异常便宜；冷启动、预热和缓存命中也要分开。否则拿热缓存结果解释所有请求，比较的起点已经不同。

基本记录包括 GPU、精度、模型版本、分辨率、帧数、采样步数、CFG 分支、并发、batch、峰值显存，以及各阶段 p50/p95。试验可以先预热三次、测二十次，但这个样本量只适合初步比较，不能保证生产尾延迟已经稳定。

\section{第一步是消除确实不变的重复工作}\label{seedance-m10-ux7b2cux4e00ux6b65ux662fux6d88ux9664ux786eux5b9eux4e0dux53d8ux7684ux91cdux590dux5de5ux4f5c}

同一提示、参考图片和固定预处理的编码结果可以缓存，键中需要包括内容哈希、裁剪缩放、编码器与适配器版本。

这里缓存的是固定输入的编码结果，和历史 K/V、跨去噪步的近似特征复用不同。后两者依赖隐藏状态，仅看文件是否变化还不够。每一种都先与关闭缓存的结果比较，确认复用的是同一个计算量，再讨论收益。

\section{注意力优化主要解决哪一类成本}\label{seedance-m10-ux6ce8ux610fux529bux4f18ux5316ux4e3bux8981ux89e3ux51b3ux54eaux4e00ux7c7bux6210ux672c}

Flash attention 改的是计算组织与显存访问，因此可以少保存中间量；完整 attention 的二次连接并没有因此消失。窗口或稀疏结构则直接改变可读范围，才减少连接数，同时也可能丢掉远距离信息。两类加速需要分别算。

GQA 主要缩小 K/V 投影与缓存，参数布局也随之改变，所以不能当作任意 MHA checkpoint 的无损选项。量化、融合与并行同样要注明作用对象，例如通信使用 FP8，不表示整个模型已经按 FP8 运行。

\section{步数减少为什么需要专门训练}\label{seedance-m10-ux6b65ux6570ux51cfux5c11ux4e3aux4ec0ux4e48ux9700ux8981ux4e13ux95e8ux8badux7ec3}

再看网络调用次数。未蒸馏模型从五十步直接减到四步，是用更粗的步长近似原来的场；另训一个四步 student，则是在重新学习少步映射。前一种做法没有自动得到后一种能力，步长增大后也更容易显露场的预测误差与积分误差。

蒸馏可以让 student 学习 teacher 的局部轨迹、跨时间一致性或预测分布；对抗与偏好监督还可用于改善少步数下的感知质量。TSCD、RayFlow 和 APT 等路线分别处理其中不同问题，不能将整套系统收益归结为一个模块。

因此蒸馏前先固定 teacher、噪声约定、scheduler、CFG 与任务分布，再同时比较速度和高运动、手部、多参考、低光、长镜头的表现。teacher 已经给错身份或动作时，student 只是更快，不能由速度推出它会纠错。

\section{VAE、量化和并行怎样选择}\label{seedance-m10-vaeux91cfux5316ux548cux5e76ux884cux600eux6837ux9009ux62e9}

VAE 靠近像素端的阶段可能占用较多时间和显存，可以考虑更窄解码器、时空分块或分辨率级联。分块需要历史特征或空间 overlap，必须检查接缝、色差和快速运动。

量化应覆盖不同噪声区间和条件数量，尤其观察低噪末段的细节、细字和音频杂声。权重量化、激活量化和缓存量化分别测量，不将一个部分的误差结论推广到全部模块。

多卡可以增加容量，也可能降低延迟，但需要付通信成本；CPU offload 少占显存，又可能受 PCIe 限制。判断是否划算，要看具体设备与并发方式。模型能装进去，和一次请求能更快完成，并不是同一个结论。

\section{各模块加速倍数为什么不能直接相乘}\label{seedance-m10-ux5404ux6a21ux5757ux52a0ux901fux500dux6570ux4e3aux4ec0ux4e48ux4e0dux80fdux76f4ux63a5ux76f8ux4e58}

设被优化部分原本占总耗时比例 \(f\)，该部分加速 \(s\) 倍，则整体加速比为：

\[
S=\frac{1}{(1-f)+f/s}.
\]

例如 DiT 占总耗时 70\%，即便它加速四倍，整体也只有：

\[
S=\frac{1}{0.3+0.7/4}\approx2.11.
\]

DiT 加速后，原来较小的 VAE 与输入处理就会占更大比例。所以下一个加速倍数的分母已经变了，不同模块、不同基线的数字不能任意相乘。网络求值次数、采样时间、吞吐和端到端延迟，都要保留各自的定义。

\section{用质量预算决定是否保留优化}\label{seedance-m10-ux7528ux8d28ux91cfux9884ux7b97ux51b3ux5b9aux662fux5426ux4fddux7559ux4f18ux5316}

每次只改变一个主要因素，固定任务、候选预算和可复用的随机条件，再做配对比较。先约定允许的质量变化与严重失败上限，避免看到结果后再放宽标准。

最后还要把结果带回交付。少步数省了十秒，若可用率下降，需要反复重生和额外审片，单位交付成本仍可能上升。加速方案的收益要连同这些新增工作一起算，不能只停在一次模型前向。

\section{真实工程从哪里进入}\label{seedance-m10-ux771fux5b9eux5de5ux7a0bux4eceux54eaux91ccux8fdbux5165}

基础采样与解码流程看 \href{https://github.com/Wan-Video/Wan2.1/blob/9737cba9c1c3c4d04b33fcad41c111989865d315/wan/text2video.py}{Wan text2video}，现成推理调用看 \href{https://github.com/modelscope/DiffSynth-Studio/blob/d32f72424cd22bc0ea8100cc57559149b4eac933/examples/wanvideo/model_inference/Wan2.1-T2V-1.3B.py}{DiffSynth 样例}，因果缓存看 \href{https://github.com/NVlabs/LongLive/blob/31c70d9071cd7851d82939806dc9fac645fcabe5/LongLive1.0/pipeline/causal_inference.py}{LongLive1.0}。先选择一个完整基线，再按具体瓶颈阅读对应模块；这些入口不意味着量化、稀疏和缓存可以在任意 checkpoint 上直接组合。
